    \begin{subsection}{Lattices}
    In this section we will provide the preliminaries from the theory of lattices for the sheaf construction
    on \( \mathcal{X} = L^\infty(\Omega, \mathcal{F}, \mathbb{P}) \), such as their definition, properties
    such as completeness and distributivity, which allow for a well-defined theory of ideals. Further a impression 
    of the geometric correspondence to acceptance sets is given, whose extension is explored in the subsequent sections.
    For the full proofs the reader may refer to the thesis \cite{mymaster}, which was the initial motivation of this work. 

    \begin{definition}
        A semilattice \( (S, \otimes ) \) is a set \( S \) with a binary operation which is 
        commutative, associative, and idempotent. It is called bounded if there's an element 
        \( e \in S \) s.t. \( x \otimes e = x \) for all \( x \in S \).
        Given two semilattices \( (L, \wedge), (L, \vee) \) on a set \( S \), whose operations
        satisfy the absorption laws:
        \begin{align*}
            a \vee (a \wedge b) = a  \\
            a \wedge (a \vee b) = a,
        \end{align*}
        the resulting structure \( (L, \wedge, \vee) \) is called a lattice and the operations \( \wedge \), \( \vee \) join and meet respectively.
        If both semilattices are bounded, the corresponding elements are denoted by \( \bot,\top \) respectively.
        A lattice is distributive if it in addition satisfies
        \begin{align*}
         a \vee (b \wedge c) = (a \wedge b) \vee (a \wedge c)  \\
         a \wedge (b \vee c) = (a \vee b) \wedge (a \vee c)  \\
        \end{align*}
    \end{definition}

    \begin{lemma}
        In a lattice \( (L, \wedge, \vee) \) for an element \( a \in L \) the sets
        \begin{align*}
            (a)^\uparrow &= \{ b \in L | a \leq b\}   \quad (a)^\downarrow = \{ b \in L | b \leq a\}  \\
            (a)^{\upuparr} &= \{ b \in L | a < b\} \quad (a)^{\dodoarr} = \{ b \in L | b < a\}  
        \end{align*}
        are semilattices, which we call the (strict) upward and (strict) downward of "\( a \)" respectively.
        If \( L \) is bounded, so are the semilattices.
    \end{lemma}

    \begin{proposition}
        There exists a unique partial order on a semilattice and for a lattice these orders agree.
    \end{proposition}
    \begin{proof}
        Define \( a \leq b \iff a \vee b = b \) and for lattices additionally \( a \leq b \iff a \wedge b = a \), 
        the rest is straightforward. 
    \end{proof}

    \begin{definition}
        A lattice is complete if each subset \( S \subset L \) has a greatest lower bound and least upper bound in \( L \),
        given by \( \bigvee_{s \in S} s \), \( \bigwedge_{s \in S} s \) respectively.
    \end{definition}
    It might appear that when a lattice is bounded, it is also complete. The two definitions coincide
    for finite \( L \), the meet and join of a lattice are only defined on pairs of elements and with that on finite \( S \), 
    while in general \( S \) might be an uncountable set. Similar distinctions should be noted for distributivity, where the counterpart 
    is completely distributive.

    \begin{definition}
        A complete distributive lattice \( L \) is join infinite distributive (JID) if it holds for any \( S \subset L \) and \( a \in L \)
        \[
            a \wedge \bigvee_{b \in S} b = \bigvee_{b \in S} a \wedge b,
        \]
        analogously \( L \) is meet infinite distributive (MID) if it holds for such \( S \) and \( a \):
        \[
            a \vee \bigwedge_{b \in S} b = \bigwedge_{b \in S} a \vee b.
        \]
    \end{definition}

    While the random variables \( \mathcal{X} : \Omega \to \mathbb{R} \) formed a ring, risk measures importantly do not with respect to 
    the same operations, e.g. \( \rho_1(X + m) + \rho_2(X + m) = \rho_1(X) + \rho_2(X) - 2m \), so cash-invariance is not maintained. 
    If risk measures are supposed to have a relation to random variables like \( \mathbb{K}[\mathbf{x}] \) has to \( \aff{\mathbb{K}}{n} \),
    as in Example \ref{algstructuresheaf}, we need a different structure, and lattices are a promising candidate.
    \begin{theorem} \label{monrmarelattice}
        The space of monetary risk measures \( \mathcal{R} \), that is functionals \( \rho : \mathcal{X} \to \mathbb{R} \), which satisfy
        the monotonicity and cash invariance from Definition \ref{defriskmeasure}, are a distributive unbounded lattice under the operations 
        \begin{align*}
            (\rho_1 \vee \rho_2)(X) &:= \max \{\rho_1(X), \rho_2(X) \} \\
            (\rho_1 \wedge \rho_2)(X) &:= \min \{\rho_1(X), \rho_2(X) \}.
        \end{align*}
        By modifying the codomain of risk measures to the extended real numbers \( \extreals = \mathbb{R} \cup \{ -\infty, \infty \} \)
        and adding the risk measures \( \rho_{\bot}(X) = -\infty,\ \rho_{\top}(X) = \infty \) to \( \mathcal{R} \) one obtains a distributive,
        bounded, complete lattice \( \extrm \) of which \( \mathcal{R} \) is a sublattice.
    \end{theorem}
    \begin{remark} \label{completionofriskmeasures}
        There is a nice, intuitive interpretation of these measures, but it comes at the cost of Lipschitz continuity, which is not too much of a burden, 
        since we have to keep only two special cases in mind. Otherwise, we do not allow \( \rho \in \extrm \) to take on the value \( - \infty \).
    \end{remark}

    \begin{lemma} \label{extrmisjidandmid}
       The lattice \( \extrm \) is JID and MID.
    \end{lemma}
    \begin{proof}
        \( \extrm \) is complete by Lemma \ref{extrmiscomplete} and the equality can be verified pointwise.
        We begin with JID. Let \( I \subseteq \extrm \), \( \rho \in \extrm \), \( X \in \mathcal{X} \)
        and let \( p = \rho(X) \) and \( q = \bigvee_{\nu \in I} \nu(X) = \bigvee_{\nu \in I} q_\nu = \sup_{\nu \in I} q_\nu \)
        be distinct members of \( \extreals \), since otherwise the result is trivial. Keep in mind that \( \extreals \) is totally ordered.
        Proceed by cases:
        \begin{itemize}
            \item \( p \geq q \): By assumption \( p \wedge q = q \) further for each \( \nu \in I \) we also have \( q_{\nu} \wedge p = q_{\nu} \)
                as otherwise we would have a contradiction. Therefore \( \bigvee_{\nu \in I} q_\nu \wedge p = \bigvee_{\nu \in I} q_\nu = q\).
            \item \( p < q \): By assumption \( p \wedge q = p \), suppose that there's no \( \nu \in I \) s.t. \( q_{\nu} > p \),
                then \( p \) is an upper bound for the set \( \{ q_{\nu} \}_{\nu \in I} \) and we would have \( \sup_{\nu \in I} q_{\nu} \leq p \)
                contradicting our hypothesis. Therefore we have some \( \nu^\ast \in I \) with \( q_{\nu^\ast} > p \) and \( \bigvee q_{\nu} \wedge p \geq q_{\nu^\ast} > p \).
        \end{itemize}
        The argument then goes analogously for MID.
    \end{proof}

    To give a taste of the intuitiveness of our construction and why lattices are a natural analog for algebraic geometry, let us turn to the geometric side:
    \begin{corollary} \label{acceptanceoperations}
        The operations in \( (\extrm, \vee, \wedge) \) correspond to intersections and unions of acceptance sets
        and vice-versa. Further the lattice order translates into acceptance sets as: \( \rho_1 \leq \rho_2 \) for \( \rho_1, \rho_2 \in \extrm \) if and only if \( \mathcal{A}_{\rho_2} \subseteq \mathcal{A}_{\rho_1} \).
    \end{corollary}
    \iffalse
    As C.G. Jung already said, "Life calls not for perfection, but for completeness" we shall 
    focus our study on \( \extrm \) from hereon as complete lattices are better behaved 
    regarding algebraic analysis. A hint at this is Definition \ref{atomisticlattice}, and as we will see, the definition or verification of ideals becomes easier if we grant ourselves this freedom. 
    \fi
\end{subsection}
